\section{Limitations and Future Work}
\label{sec:limitations}

\subsection{Computational Complexity}

Our model sets no limit on the amount of computation a strategy can perform
between steps of the universe, and our proposed potential-maximizing strategy
requires vast computational resources to search through a large space of
actions in a complex universe. This is clearly unrealistic.

Nevertheless, we believe the model is still useful. In practice, when problems
become computationally infeasible, it is often possible to find heuristics or
approximations that perform well for commonly encountered problems. A formal
model of what an ideal strategy would be helps guide the development and
evaluation of such approximations.

\subsection{All-Knowingness}

Our existence proof for potential-maximizing strategies assumes that the
strategy has full knowledge of the universe's state, the transition function,
the reward function, and all the other strategies. While strategies do not
always have such vast knowledge to begin with, they can often learn at least
approximations of it.

For example, formalisms like AIXI~\citep{Hutter:04uaibook} show how strategies
can both maximize their reward and learn the workings of their universe (and the
actors within it), using methods like Occam's razor. AIXI is incomputable, but
there are computable approximations of it. Even then, the discussion on
computability from the previous section would still apply.

\subsection{Fixed Points}

We have shown that if the universe, the reward function, and the other players'
strategies are known, a potential-maximizing strategy always exists.

This is, however, not entirely satisfactory. While it does seem reasonable that
a strategy could learn its universe, its reward function, and the other players'
strategies, the full knowledge of the other players' strategies can sometimes
give the potential-maximizing strategy an unfair advantage. Consider, for
example, a universe where agents play a game of Rock-Paper-Scissors. While
our proposed potential-maximizing strategy has to win against every strategy, it
also knows exactly what strategy it plays against and can thus always respond
with the winning action (Scissors against Paper, etc.). This problem may be
avoided by using a fixed-point theorem, similar to the approach taken in Nash
Equilibria.

We believe our model still provides novel and interesting insights, but an
extension of the model to incorporate fixed-point theorems seems like an
important avenue for future work.

\subsection{Randomness}

Our model is completely deterministic. Lifting this assumption would help to
model the physical universe more accurately (e.g., quantum physics is inherently
non-deterministic) and may also help with the incorporation of fixed-point theorems by
allowing mixed strategies.

\subsection{Discrete Time}

Our model assumes that time progresses in discrete steps. This seems like a
reasonable assumption, as the step length can be arbitrarily small.

% \subsection{Formalizing the Discussion}

% This paper's \cref{sec:discussion} is a mostly informal discussion. However,
% some of the concepts discussed there, like the substitutability of reward
% functions, appear amenable to formal modeling. This seems like an interesting
% avenue for future work.
